\scanpage{193}
\mspage{3}
\margin{Utilise R1 et R2}
\uline{Prop 2} Soit $X\in\mathrm{Ob}\,\mathscr R_{S'}$. Alors $X=\varnothing\Longleftrightarrow r(X)=\varnothing$, $X$ connexe $\Longleftrightarrow r(X)$ connexe.

Soit $C$ une \inserted{$\mathcal U$-}\deleted{\unsure{cat.}}catégorie multigaloisienne, $P$ un objet.
% "catégorie" is crossed but has a dotted reinstatement.
Un objet $x$ de $C$ est dit \uline{subordonné} à $P$ si $X_P=\deleted{\unsure{P\times X}}X\times P\longrightarrow P$ est \deleted{\unsure{loc. constant pour la top}}\inserted{tel que, sur le ss-objet $P_n$ de $P$ où $X_P$ est de rang $n$, \deleted{$X$} est \uline{constant}.}
\diagram{\begin{tikzcd}[ampersand replacement=\&]
X\arrow[d,dash] \& X\times P\arrow[l,dash]\arrow[d,dash]\\
e \& P\arrow[l,dash]
\end{tikzcd}}
Pour $P$ fixé, les objets $X$ de $C$ subordonnés à $P$ forment une ss-catégorie pleine stable par $\mathop{\mathrm{lim}}\limits_{\substack{\longleftarrow\\\longrightarrow}}$ finies et par sommande direct [car les objets \inserted{sur $P$} qui sont ``constants par morceaux'' forment une \inserted{ss-}catégorie pleine de $C_{/P}$ stable par $\mathop{\mathrm{lim}}\limits_{\substack{\longrightarrow\\\longleftarrow}}$ finies et par sommandes directes \ldots],
\margin{dégager en lemme de catégorie multigaloisienne}
% slanted
donc c'est une ss-catégorie multigaloisienne pleine de $C$, soit $C^{(P)}$. Si $P$ est galoisien de groupe $G$, cette catégorie est équivalente à celle des objets \inserted{dans $P$} cts par morceaux, \deleted{\unsure{munis}} avec op. de $G$ dessus compatible avec son op. sur $P$.

Si maintenant $C=\mathscr R_S$, on peut regarder la ss-catégorie fibrée pleine de $\mathscr R$ formée des $X'\in\mathrm{Ob}\,\mathscr R_{S'}$ qui sont \uuline{localement} \uline{subordonnés} à $P$. On voit encore que c'est stable par $\mathop{\mathrm{lim}}\limits_{\substack{\longrightarrow\\\longleftarrow}}$ et sommandes directes, donc c'est un ss-\deleted{catégorie}\inserted{champ en} \deleted{fibrée}\inserted{catégories multigaloisiennes} plein $\mathscr R^P$ de $\mathscr R$.

Supposons maintenant $P$ galoisien de groupe $G$, alors les objets de $\mathscr R_S$ \uline{subordonnés} (pas seul$^{\rm t}$ loc$^{\rm t}$ !) à $P$ correspondent aux objets sur $P$ \uline{constants} par morceaux, avec op. de $G$ dessus. Or, grâce à R1, R\unsure{4|3|\deleted{2}4}, un objet \inserted{$X$} sur $P$ constant par morceaux définit un objet \inserted{$r(X)$} sur $r(P)$ constant par morceaux, et
\margin{schémas}
% archive sheet number: 192
